\section{Conclusion}
\label{sec:conclusion}

We evaluated record-level supervisory analytics for SOCs as implemented in SAT-SA:
rule-based analysis of the records a monitored organization submits, saturating risk
fusion and ranking, a human supervisory decision, and cryptographic binding of that
decision to the analysed state. Under controlled conditions the mechanisms behave as
specified---declared findings are emitted, validation follows its contract, the tested
alterations of the decided record are detected, and the evaluated interruptions
recover---but simple rules match much of the detection and ranking performance even on
generated populations that favour SAT-SA by construction, not every analytical worker
improved the ranking metric, and stale and contradictory evidence passes silently. On a public IT incident log the
supervisory risk did not track the only outcome available, because that outcome
measures timeliness, which the fast-closure signal treats in the opposite direction,
and because existence rules saturate at scale. What can be concluded is therefore
mechanism-level: the pipeline is inspectable and behaves as designed, while its
usefulness for supervision is untested. The evidence needed next is not more
controlled experiments but SOC records with independent outcomes held out from
development, expert-labelled findings, a study with supervisors, prevalence-aware
scoring and consistency checks for stale and conflicting records, and validation of
the containerized PostgreSQL deployment.
